\documentclass[11pt,a4paper]{article}
\usepackage[T1]{fontenc}
\usepackage[utf8]{inputenc}
\usepackage{lmodern}
\usepackage[margin=27mm]{geometry}
\usepackage{amsmath,amssymb,amsthm,mathtools}
\usepackage{microtype,booktabs,longtable,array,enumitem,xcolor}
\usepackage[colorlinks=true,linkcolor=blue!45!black,citecolor=blue!45!black,urlcolor=blue!50!black,breaklinks=true]{hyperref}
\usepackage{fancyhdr}
\pagestyle{fancy}\fancyhf{}\fancyhead[L]{\small Keio / Research manuscript}\fancyhead[R]{\small September 2026}\fancyfoot[C]{\thepage}
\setlength{\headheight}{14pt}
\setlength{\emergencystretch}{3em}
\setlength{\parskip}{0.35em}
\setlist{itemsep=2pt,topsep=4pt}
\newtheorem{theorem}{Theorem}[section]
\newtheorem{lemma}[theorem]{Lemma}
\newtheorem{proposition}[theorem]{Proposition}
\newtheorem{corollary}[theorem]{Corollary}
\newtheorem{conjecture}[theorem]{Conjecture}
\newtheorem{criterion}[theorem]{Criterion}
\theoremstyle{definition}\newtheorem{definition}[theorem]{Definition}
\newtheorem{example}[theorem]{Example}
\theoremstyle{remark}\newtheorem{remark}[theorem]{Remark}
\newtheorem{audit}[theorem]{Proof-status note}
\newcommand{\Q}{\mathbb Q}\newcommand{\Z}{\mathbb Z}\newcommand{\N}{\mathbb N}\newcommand{\R}{\mathbb R}\newcommand{\C}{\mathbb C}
\newcommand{\Qbar}{\overline{\mathbb Q}}
\DeclareMathOperator{\Tr}{Tr}\DeclareMathOperator{\Norm}{N}\DeclareMathOperator{\trdeg}{trdeg}\DeclareMathOperator{\ord}{ord}\DeclareMathOperator{\lcm}{lcm}
\DeclareMathOperator{\Ein}{Ein}\DeclareMathOperator{\Ei}{Ei}\DeclareMathOperator{\Ci}{Ci}\DeclareMathOperator{\Chi}{Chi}\DeclareMathOperator{\Li}{Li}
\newcommand{\floor}[1]{\left\lfloor #1\right\rfloor}\newcommand{\ceil}[1]{\left\lceil #1\right\rceil}
\newcommand{\sourceT}[1]{\textup{[source T#1]}}
\author{Keio}
\date{5 September 2026\\\small Standalone research draft; source audit identifiers retained}
\title{Additive Clocks, Non-Reversible Recurrences,\\ and Transcendental Prime-Representing Constants}
\newcommand{\PP}{\mathcal P}
\newcommand{\WW}{\mathcal W}
\newcommand{\dens}{\mathrm d}
\begin{document}
\maketitle
\begin{abstract}
For integers $a\ge2$, $b\ge1$ and $\varepsilon\in\{-1,1\}$, the
sequence $ba^n+\varepsilon n$ visits every residue class modulo every
positive integer with positive natural density. We prove this by
induction on the modulus, lowering an eventual period to its greatest
common divisor with the modulus. The sequences admit an inhomogeneous
recurrence of order two but no reversible integer recurrence, even
after changing the order. Pulling back a theorem of Dubickas gives
positive lower density of simultaneous composite values in each fixed
integer window about $\lfloor\alpha^{ba^n+\varepsilon n}\rfloor$ for
every Pisot or Salem number $\alpha$. This supplies sequences requested
in Saito's Problem~1.7. A separate application of Saito's Type-B
criterion proves transcendence of the least prime-representing
constants for $ba^k\pm k$ when $a\ge3$, and for two carefully indexed
alternating-sign sequences. For the negative-sign family, we also
prove an exponentially sharp asymptotic about the algebraic anchor
$2^{1/(a-1)}$. The anchor itself fails infinitely many prime conditions.
The manuscript distinguishes proved applications from the external
transcendence criterion and makes no claim to resolve Mills'
constant or to replace prime search by a numerical approximation.
\end{abstract}

\section{The questions and the organizing principle}
The non-reversible recurrence problem asks for an integer sequence
$R(n)$, governed by a linear recurrence with a constant inhomogeneous
term, such that
\[
 \floor{\alpha^{R(n)}}
\]
is composite infinitely often for every Pisot number $\alpha$.
This is Problem~1.7 in Saito's preprint \cite{SaitoIntervals}. The
sequence is to be chosen independently of $\alpha$. Non-reversibility
means that the trailing recurrence coefficient cannot be $1$ or
$-1$; below we prove the stronger fact that no other recurrence
representation makes our sequences reversible.

Our examples are
\begin{equation}\label{eq:clock}
 R_{a,b,\varepsilon}(n)=ba^n+\varepsilon n,
 \qquad a\ge2,\ b\ge1,\ \varepsilon\in\{-1,1\}.
\end{equation}
The same elementary congruence property has another application.
For a positive integer sequence $C=(C_k)_{k\ge1}$, define
\begin{equation}\label{eq:Wdef}
 \WW(C)=\{x>1:\floor{x^{C_k}}\in\PP\text{ for all }k\ge1\},
 \qquad \xi(C)=\min\WW(C)
\end{equation}
whenever the minimum exists; $\PP$ is the set of positive primes.
For the exponent sequences $C_k=ba^k+\varepsilon k$, $a\ge3$, the
congruence property verifies all arithmetic assumptions of Saito's
Type-B theorem \cite{SaitoVariants}. It follows that every one of these
specified minima exists and is transcendental.

The term \emph{additive clock} refers simply to the linear term
$\varepsilon n$. When the exponential part returns modulo $q$, this
term moves by $\varepsilon t$, where $t$ is an eventual period.
Although this increment may fail to generate all residues modulo
$q$, it generates all residues compatible modulo $\gcd(t,q)$.
The latter modulus is strictly smaller. Induction completes the
missing residues.

The general congruence theorem is the shared input, not a new proof of
Saito's transcendence theorem. We give proofs of the congruence,
recurrence, pullback, braid and asymptotic statements, while explicitly
citing the two external number-theoretic inputs. Source identifiers
T47--T49 and T56--T57 refer to the author's initial research notes.
The special case T48 is absorbed into the stronger coefficient theorem.

\section{Every residue occurs with positive density}
For $S\subseteq\N$, let $\dens(S)$ denote its natural density when
that density exists, and let $\underline{\dens}(S)$ denote its lower
natural density.

\begin{theorem}[Universal residues; \sourceT{47}, \sourceT{56(a)}]
\label{thm:residues}
Fix the parameters in \eqref{eq:clock}. For every integer $q\ge1$
and every $c\in\Z$, the set
\[
 S(q,c)=\{n\ge1:R_{a,b,\varepsilon}(n)\equiv c\pmod q\}
\]
has positive natural density. It is an eventually periodic subset of
$\N$, and it contains all sufficiently large members of at least one
arithmetic progression.
\end{theorem}

\begin{lemma}[A short eventual period]\label{lem:period}
Let $a\ge2$ and $q>1$. There exist positive integers $e,t$ such that
\begin{equation}\label{eq:period}
 a^{n+t}\equiv a^n\pmod q\quad(n\ge e),\qquad 1\le t<q.
\end{equation}
\end{lemma}
\begin{proof}
Write $q=uv$, where every prime divisor of $u$ divides $a$ and
$\gcd(a,v)=1$. More precisely, $u$ is the product of the prime-power
factors of $q$ supported on primes dividing $a$. Thus $\gcd(u,v)=1$.
For all sufficiently large $n$, $u\mid a^n$. If $v>1$, choose
$t=\ord_v(a)$. Then
\[
 t\le\varphi(v)<v\le q,
\]
and \eqref{eq:period} follows from the Chinese remainder theorem.
If $v=1$, choose $t=1<q$, since the exponential part is eventually
zero modulo $q$.
\end{proof}

\begin{proof}[Proof of Theorem~\ref{thm:residues}]
We first prove the existence of arbitrarily large solutions, using
strong induction on $q$. For $q=1$ every index works. Suppose $q>1$,
and take $e,t$ as in Lemma~\ref{lem:period}. Set
\[
 g=\gcd(t,q)<q.
\]
By the induction hypothesis, there are arbitrarily large $n\ge e$
with $R_{a,b,\varepsilon}(n)\equiv c\pmod g$. For such an $n$, the
linear congruence
\begin{equation}\label{eq:jcong}
 \varepsilon jt\equiv c-R_{a,b,\varepsilon}(n)\pmod q
\end{equation}
is solvable because its right-hand side is divisible by $g$.
Choose a nonnegative solution $j$ and put $N=n+jt$. The period
identity gives $a^N\equiv a^n\pmod q$, so
\[
 R_{a,b,\varepsilon}(N)
 \equiv R_{a,b,\varepsilon}(n)+\varepsilon jt
 \equiv c\pmod q.
\]
Since $N\ge n$, solutions can be chosen arbitrarily large. The
coefficient $b$ is unchanged throughout and imposes no coprimality
condition.

Now put $Q=\lcm(q,t)$. For every $n\ge e$,
\[
 R_{a,b,\varepsilon}(n+Q)
 \equiv R_{a,b,\varepsilon}(n)\pmod q.
\]
Thus $S(q,c)$, apart from finitely many indices, is a union of residue
classes modulo $Q$. At least one such class is present by the first
part of the proof. Consequently its density exists and is at least
$1/Q>0$.
\end{proof}

\begin{remark}[Surjectivity is not uniform distribution]
The theorem supplies at least one eventual index class for each target
residue. It does not assert that the $q$ residues have equal density.
The period $Q$ may depend on $a$ and $q$, and different residues can
have different numbers of representatives modulo $Q$.
\end{remark}

\section{Intrinsic non-reversibility}
An integer sequence $R$ is an inhomogeneous linear recurrence sequence
of order $d$ if, for fixed integers $a_0,\ldots,a_{d-1},c$ with
$a_0\ne0$,
\[
 R(n+d)=a_{d-1}R(n+d-1)+\cdots+a_0R(n)+c.
\]
Such a representation is reversible when $a_0=\pm1$. In a sequence
with several possible recurrence representations, merely displaying
one with $a_0\ne\pm1$ does not exclude a hidden reversible one. We
therefore check all integer representations at once.

\begin{proposition}\label{prop:nonrev}
The sequence in \eqref{eq:clock} is positive and strictly increasing
for $n\ge1$. It satisfies the order-two relation
\begin{equation}\label{eq:recurrence}
 R(n+2)=(a+1)R(n+1)-aR(n)+\varepsilon(1-a).
\end{equation}
Every integer inhomogeneous recurrence representation of this
sequence has trailing coefficient divisible by $a$. In particular,
none is reversible.
\end{proposition}
\begin{proof}
The exponential and linear terms directly verify
\eqref{eq:recurrence}. Positivity follows from
$ba^n-n\ge2^n-n\ge1$. Also
\[
 R(n+1)-R(n)=b(a-1)a^n+\varepsilon\ge1.
\]
Let $E$ be the shift operator, and suppose $Q(E)R$ is constant for
some monic integer polynomial $Q$ associated to an inhomogeneous
recurrence. The exponential term in $Q(E)R$ is $ba^nQ(a)$;
the coefficient of $n$ is $\varepsilon Q(1)$. An exponential plus
a polynomial can be constant for all positive $n$ only if
\[
 Q(a)=0,\qquad Q(1)=0.
\]
Since $a\ne1$, the monic polynomial $(X-a)(X-1)$ divides $Q$ in
$\Z[X]$. Thus $a\mid Q(0)$. The trailing coefficient of the
recurrence is $-Q(0)$ and therefore cannot be $1$ or $-1$.
This also shows directly that the homogeneous minimal annihilator is
$(X-a)(X-1)^2$.
\end{proof}

The term ``order two'' in Proposition~\ref{prop:nonrev} refers to a
recurrence with a constant inhomogeneous term. Its homogeneous
annihilator has degree three. Conflating these two conventions would
misstate the non-reversibility conclusion.

\section{Pulling back composite progressions}
We use the following established result as an external input. It is
Theorem~1.1 of Dubickas \cite{Dubickas}, also recorded as (D1) in
\cite{SaitoIntervals}.

\begin{theorem}[Dubickas, external input]\label{thm:Dub}
Let $(T_m)_{m\ge1}$ be an unbounded integer linear recurrence sequence.
For each integer $J\ge1$, there are positive integers $L,m_0$ and
primes $p_j$, $|j|\le J$, such that
\[
 p_j\mid T_{Lk+m_0}+j\qquad(k\ge1,\ |j|\le J).
\]
The primes $p_j$ need not be distinct.
\end{theorem}

We pass from trace sequences to floors with a bounded-error argument.
This step only requires a finite collection of available divisors.
It does not require one particular prime to divide every floor value
with the same external shift.

\begin{lemma}[Composite floor windows]\label{lem:floorAP}
Let $\alpha$ be a Pisot or Salem number, and fix $H\ge0$. There
exist positive integers $L,m_0,K$ such that
\begin{equation}\label{eq:floorAP}
 \floor{\alpha^{Lk+m_0}}+h
 \quad\text{is composite for all }k\ge K,\quad |h|\le H.
\end{equation}
\end{lemma}
\begin{proof}
Let $d$ be the degree of $\alpha$ and let
$\alpha=\alpha_1,\ldots,\alpha_d$ be its conjugates. The integer
trace sequence
\[
 T_m=\sum_{i=1}^d\alpha_i^m
\]
satisfies the linear recurrence given by the monic minimal polynomial.
For a Pisot number the other conjugates have modulus less than $1$;
for a Salem number they have modulus at most $1$. Hence
\[
 |T_m-\alpha^m|\le d-1,\qquad
 e_m:=\floor{\alpha^m}-T_m\in\Z,\quad |e_m|\le d.
\]
The same bound shows $T_m\to\infty$, so the trace recurrence is
unbounded. Apply Theorem~\ref{thm:Dub} with $J=H+d$. Along the
resulting progression, every $T_{Lk+m_0}+j$, $|j|\le J$, has the
fixed divisor $p_j$. For sufficiently large $k$, all these positive
integers exceed all the finitely many primes $p_j$. Finally,
\[
 \floor{\alpha^{Lk+m_0}}+h
 =T_{Lk+m_0}+(e_{Lk+m_0}+h),
\]
and $|e_{Lk+m_0}+h|\le J$. Choosing the available prime for this
integer shift proves proper divisibility and \eqref{eq:floorAP}.
\end{proof}

\begin{theorem}[A solution of the non-reversible construction problem;
\sourceT{49}]\label{thm:composite}
For every $a\ge2$, $b\ge1$, $\varepsilon=\pm1$, every Pisot or
Salem number $\alpha$, and every integer $H\ge0$, the set
\[
 \left\{n\ge1:
 \floor{\alpha^{ba^n+\varepsilon n}}+h\text{ is composite for all }
 h\in\Z,\ |h|\le H\right\}
\]
has positive lower natural density.
\end{theorem}
\begin{proof}
Choose $L,m_0,K$ from Lemma~\ref{lem:floorAP}. By
Theorem~\ref{thm:residues}, the set
\[
 S=\{n\ge1:R_{a,b,\varepsilon}(n)\equiv m_0\pmod L\}
\]
has positive natural density. Since $R_{a,b,\varepsilon}(n)\to\infty$,
all sufficiently large $n\in S$ satisfy
$R_{a,b,\varepsilon}(n)=Lk_n+m_0$ with $k_n\ge K$.
Equation~\eqref{eq:floorAP} applies to each of them. Removing finitely
many indices from $S$ leaves its density unchanged, proving the
claimed lower density.
\end{proof}

Taking $b=1$, $H=0$ and restricting to Pisot numbers gives the exact
construction asked for in Saito's Problem~1.7. The clock is independent
of $\alpha$, and Proposition~\ref{prop:nonrev} verifies the required
non-reversibility. The Salem and fixed-window statements are further
applications of the same proof. The theorem does not replace the
exponent $3^r$ in the separate cubic Pisot problem by an equivalent
one: $a^r\pm r$ is a different sequence.

\section{The external Type-B criterion, with all hypotheses}
The following is the transcendence conclusion of Theorem~2.3 in
Saito's version~3 \cite{SaitoVariants}, expressed directly in terms of
the exponent sequence. This is an external theorem, not a lemma
proved in this manuscript.

\begin{criterion}[Saito Type B]\label{crit:typeB}
Let $C_k\in\N$, set $c_1=C_1$ and $c_{k+1}=C_{k+1}/C_k$, and assume
\begin{align*}
 &\textup{(B1)}\quad c_1\ge1,\qquad
 \textup{(B2)}\quad c_{k+1}\ge2\ (k\ge1),\\
 &\textup{(B3)}\quad\limsup_{k\to\infty}c_{k+1}>40/19,\qquad
 \textup{(B4)}\quad C_k=c_1\cdots c_k\in\N.
\end{align*}
Choose sufficiently small $\eta>0$ so that
$I=\{k:c_{k+1}\ge40/19+\eta\}$ is infinite. Assume also
\begin{align*}
 \textup{(B5)}\quad&\forall m\ge1\ \exists k\in I,
 \quad k>m,\quad C_m\mid C_k;\\
 \textup{(B6)}\quad&\forall m,L\ge1\ \exists k\in I,
 \quad k>m,\quad C_k\equiv C_m\pmod{LC_m}.
\end{align*}
Then $\xi(C)$ exists and is transcendental.
\end{criterion}

The existence of a minimum causes no extra difficulty once the
prime-representing set is nonempty. We include the basic order
argument, which is also useful for understanding the exact endpoints
of the intervals used later.

\begin{lemma}[Existence of a minimum from nonemptiness]\label{lem:min}
If $C_k$ are positive integers and $\WW(C)\ne\varnothing$, then
$\WW(C)$ has a minimum.
\end{lemma}
\begin{proof}
Let $x=\inf\WW(C)$. The first prime condition gives
$x\ge2^{1/C_1}>1$. There is a sequence $x_j\in\WW(C)$ with
$x_j\ge x$ and $x_j\to x$. For each fixed $k$, the values
$x_j^{C_k}$ approach $x^{C_k}$ from above. For all sufficiently
large $j$, their floors equal $\floor{x^{C_k}}$, including when
$x^{C_k}$ itself is an integer. This eventually constant floor is
prime because every $x_j\in\WW(C)$. Thus $x\in\WW(C)$.
\end{proof}

Criterion~\ref{crit:typeB} is used only after checking every one of its
hypotheses. In particular, a congruence result at arbitrarily large
indices does not automatically place those indices in $I$. Here we
will have $I$ cofinite, making the required intersection explicit.

\section{Transcendence for all coefficients and both signs}
\begin{theorem}[Coefficient family; \sourceT{48}, \sourceT{56(a)}]
\label{thm:family}
For every integer $a\ge3$, every integer $b\ge1$, and both choices
of $\varepsilon\in\{-1,1\}$, the constant
\[
 \xi\bigl((ba^k+\varepsilon k)_{k\ge1}\bigr)
\]
exists and is transcendental.
\end{theorem}
\begin{proof}
Put $C_k=ba^k+\varepsilon k$. Then $C_1=ba+\varepsilon\ge2$,
and every $C_k$ is a positive integer. Moreover,
\begin{equation}\label{eq:ratioineq}
 C_{k+1}-2C_k=b(a-2)a^k+\varepsilon(1-k)>0.
\end{equation}
For $\varepsilon=-1$ both terms on the right are nonnegative and the
first is positive. For $\varepsilon=1$, use
$b(a-2)a^k\ge3^k>k-1$. This verifies (B1), (B2) and (B4).
Also
\[
 \frac{C_{k+1}}{C_k}\longrightarrow a>40/19,
\]
so (B3) holds. Fix any sufficiently small
$0<\eta<a-40/19$. Then $I$ in Criterion~\ref{crit:typeB} is cofinite.

For a given $m$, apply Theorem~\ref{thm:residues} with modulus
$C_m$ and target $0$. It supplies arbitrarily large $k$ with
$C_m\mid C_k$. Choose one with $k>m$ and $k\in I$, verifying (B5).
For given $m,L$, apply the same theorem with modulus $LC_m$ and
target $C_m$. Again choose a sufficiently large solution in $I$;
this gives (B6). Criterion~\ref{crit:typeB} proves the assertion.
\end{proof}

For $a=3$, $b=1$, the first few exponents are
\[
 (3^k-k)_{k\ge1}=(2,7,24,77,238,\ldots),
\]
\[
 (3^k+k)_{k\ge1}=(4,11,30,85,248,\ldots).
\]
These are two different exponent sequences defining two minima. Their
transcendence follows from the same theorem, but this does not
establish that the minima are unequal or order them relative to
Mills' constant.

\begin{corollary}[Adjacent coefficient lines]\label{cor:adjacent}
Let $A_k=3^k-k$ and $B_k=3^k+k$. For every integer $r\ge0$, both
exponent sequences
\[
 (r+1)A_k+rB_k=(2r+1)3^k-k,
\]
\[
 rA_k+(r+1)B_k=(2r+1)3^k+k
\]
define transcendental least prime-representing constants.
\end{corollary}
\begin{proof}
Apply Theorem~\ref{thm:family} with $a=3$, $b=2r+1$ and the
corresponding sign.
\end{proof}
The coefficient theorem also covers even $b$, which are not obtained
from these adjacent nonnegative coefficient pairs. Nothing in the
congruence proof depended on $b$ being odd.

\section{Alternating signs and the necessary indexing}
Define
\begin{equation}\label{eq:braids}
 P_k=3^k+(-1)^k k,\qquad Q_k=3^k-(-1)^k k.
\end{equation}
The notation $P_k$ here is an exponent sequence, not a prime. We
spell out the starting indices because one of the sequences fails a
Type-B growth condition at its first step.

\begin{lemma}[Residues of the braids]\label{lem:braidres}
Each sequence in \eqref{eq:braids} visits every residue class modulo
every positive integer with positive natural density.
\end{lemma}
\begin{proof}
For either fixed sign $\varepsilon$, one has
$3^k+\varepsilon k\equiv1+k\pmod2$.
First suppose $q$ is even and a residue $c$ modulo $q$ is given.
The required index parity is $t\equiv c-1\pmod2$. On that parity,
$P_k=3^k+\varepsilon k$ with $\varepsilon=(-1)^t$.
Theorem~\ref{thm:residues} for this fixed sign supplies positive-density
solutions modulo $q$, and the congruence modulo $2$ forces every one
of them to have parity $t$. They are therefore solutions for $P_k$.
For $Q_k$, choose $\varepsilon=-(-1)^t$ instead.

If $q$ is odd, choose a lift $c'$ modulo $2q$ with $c'\equiv c\pmod q$
and any prescribed parity. Apply the even-modulus argument to $2q$.
This supplies a positive-density subset of the desired solutions
modulo $q$. Finally, both braid sequences are eventually periodic
modulo $q$: combine an eventual period of $3^k$ with the periods
$q$ of $k$ and $2$ of $(-1)^k$. Thus their full solution sets have
natural densities, which must be positive.
\end{proof}

\begin{theorem}[Braid constants; \sourceT{56(b)}]\label{thm:braids}
Both
\[
 \xi\bigl((P_k)_{k\ge1}\bigr)
 \quad\text{and}\quad
 \xi\bigl((Q_k)_{k\ge2}\bigr)
\]
exist and are transcendental. In the second expression the sequence
is reindexed from $1$, namely $C_j=Q_{j+1}$.
\end{theorem}
\begin{proof}
Direct subtraction gives
\begin{align*}
 P_{k+1}-2P_k&=3^k-(-1)^k(3k+1),\\
 Q_{k+1}-2Q_k&=3^k+(-1)^k(3k+1).
\end{align*}
For $k\ge2$, the inequality $3^k>3k+1$ follows by induction, so both
differences are positive. Also $P_1=2$, $P_2=11$ satisfy the first
growth condition for the $P$ sequence. For the $Q$ sequence,
$Q_1=4$, $Q_2=7$ give $Q_2<2Q_1$, explaining the stated truncation.
The reindexed sequence starts at $Q_2=7$ and has all subsequent
ratios greater than $2$. Both ratios tend to $3$.

Consequently (B1)--(B4) hold for the two stated sequences, and $I$ is
cofinite. Lemma~\ref{lem:braidres}, unaffected by deleting finitely
many indices, verifies (B5) and (B6) exactly as in the proof of
Theorem~\ref{thm:family}. Apply Criterion~\ref{crit:typeB}.
\end{proof}

The proof makes no assertion about the minimum defined by the full
$Q$ sequence starting at $k=1$. Removing a prime condition can change
the minimum, so a tail theorem must not be silently promoted to the
full sequence.

\section{Scaling, synchronization and elementary combinations}
\begin{proposition}[Exact scaling; \sourceT{56(c)}]\label{prop:scaling}
If $C$ is a positive integer sequence with $\WW(C)\ne\varnothing$
and $d\ge1$ is an integer, then
\[
 \xi\bigl((dC_k)_{k\ge1}\bigr)=\xi(C)^{1/d}.
\]
\end{proposition}
\begin{proof}
For $x>1$,
\[
 x\in\WW(dC)\quad\Longleftrightarrow\quad x^d\in\WW(C).
\]
The map $x\mapsto x^d$ is an increasing bijection of $(1,\infty)$.
It therefore carries the minimum of the first set to the minimum of
the second, which exists by Lemma~\ref{lem:min}.
\end{proof}

For $A_k=3^k-k$, $B_k=3^k+k$ and $M=\xi((3^k)_{k\ge1})$,
Proposition~\ref{prop:scaling} gives
\begin{equation}\label{eq:Millsboundary}
 \xi((A_k+B_k)_{k\ge1})=\sqrt M.
\end{equation}
This is an identity about the \emph{sum of exponent sequences}. It is
not an identity involving $\xi(A)+\xi(B)$ or $\xi(A)\xi(B)$.
The transcendence of the left side is equivalent to that of $M$;
the separate transcendence statements for $\xi(A)$ and $\xi(B)$ do
not decide it.

\begin{proposition}[A simultaneous-return obstruction;
\sourceT{56(d)}]\label{prop:sync}
For every $m\ge1$, no integer $k>m$ satisfies both
\[
 A_k\equiv A_m\pmod{3^{m+1}},\qquad
 B_k\equiv B_m\pmod{3^{m+1}}.
\]
\end{proposition}
\begin{proof}
Since $2$ is invertible modulo $3^{m+1}$, adding the two congruences
would give $3^k\equiv3^m\pmod{3^{m+1}}$. But $k\ge m+1$ makes
the left side zero and the right side nonzero modulo this modulus.
\end{proof}

The proposition does not contradict individual Type-B returns. Their
indices need not coincide. It rules out this particular attempt to
synchronize the two returns and cancel the linear terms, rather than
all possible relations between the two constants.

\begin{corollary}[Elementary combinations; \sourceT{56(e)}]
\label{cor:combinations}
Let $\alpha=\xi(A)$ and $\beta=\xi(B)$. At least two of
\[
 \alpha+\beta,\qquad \alpha\beta,\qquad \alpha/\beta
\]
are transcendental. Also $\alpha+i\beta$ is transcendental.
\end{corollary}
\begin{proof}
Both $\alpha$ and $\beta$ are positive transcendental numbers by
Theorem~\ref{thm:family}. If their sum and product were algebraic,
they would be roots of a quadratic polynomial with algebraic
coefficients. If their sum and ratio were algebraic, then
$\beta=(\alpha+\beta)/(1+\alpha/\beta)$ would be algebraic.
If their product and ratio were algebraic, then
$\alpha^2=(\alpha\beta)(\alpha/\beta)$ would be algebraic.
Every case contradicts the transcendence of $\alpha$ or $\beta$.
Finally, the real and imaginary parts of an algebraic complex number
are algebraic, excluding algebraicity of $\alpha+i\beta$.
\end{proof}
These are algebraic consequences of the two known transcendental
values. They do not imply algebraic independence.

\section{An algebraic anchor that fails infinitely many conditions}
For the rest of the manuscript let
\begin{equation}\label{eq:anchor}
 C_k=a^k-k,\quad \xi_a=\xi(C),\quad R_a=2^{1/(a-1)},
 \quad r_a=\frac{a^2-2}{a-1}=a+1-\frac1{a-1},
 \qquad a\ge3.
\end{equation}
The constant $R_a$ is the smallest number compatible with a first
prime value at the first exponent. It is a useful scale, but it does
not satisfy the full prime conditions.

\begin{proposition}[Exact anchor obstruction; \sourceT{57(a)}]
\label{prop:anchor}
For every integer $a\ge3$,
\[
 \xi_a>R_a.
\]
For every $k>1$ with $k\equiv1\pmod{a-1}$, the value
$R_a^{C_k}$ is an integer power of $2$ with exponent greater than $1$,
and hence its floor is composite.
\end{proposition}
\begin{proof}
The first prime condition gives
$\xi_a^{a-1}\ge2$, so $\xi_a\ge R_a$. Since
$a\equiv1\pmod{a-1}$,
\[
 C_k\equiv1-k\pmod{a-1}.
\]
At the stated indices $m_k=C_k/(a-1)$ is a positive integer greater
than $1$. Therefore
$R_a^{C_k}=2^{m_k}$ is composite. Such indices are infinite, so
$R_a\notin\WW(C)$ and the lower bound for its minimum is strict.
\end{proof}

This strict inequality is independent of the transcendence proof.
For example, $R_3=\sqrt2$ gives floors $2$ and $11$ at exponents
$2$ and $7$, but at the next exponent $24$ it gives $4096$. Agreement
with the first two prime conditions therefore provides no all-index
certificate.

\section{Short prime intervals and nested prime cylinders}
The analytic external input is the Baker--Harman--Pintz prime interval
theorem \cite{BHP}, in the right-interval form also stated in
\cite[Theorem~1.6]{SaitoVariants}. With
\[
 \theta=21/40,
\]
there is a constant $X_0$ such that, for every $X\ge X_0$, a prime
exists strictly between $X$ and $X+X^\theta$, with the upper endpoint
allowed. The usual quantitative version gives many primes, so
excluding a possible prime at $X$ does not alter this assertion for
sufficiently large $X$. No explicit numerical $X_0$ is used here.

\begin{lemma}[Room for the next cylinder]\label{lem:room}
For every $p\ge2$ and every real $r\ge3$,
\[
 (p+1)^r-p^r>p^{\theta r}+1.
\]
\end{lemma}
\begin{proof}
The integral form of the mean value estimate gives
$(p+1)^r-p^r>rp^{r-1}$. Since
\[
 r-1-\theta r=(19/40)r-1\ge17/40>0,
\]
one has $rp^{r-1}\ge3p^{\theta r}>p^{\theta r}+1$.
\end{proof}

\begin{lemma}[A branch beginning at $2$]\label{lem:branch}
There exists an integer $a_0$ such that, for every integer $a\ge a_0$,
a number $A_a\in\WW((a^k-k)_{k\ge1})$ can be chosen with
\[
 \floor{A_a^{a-1}}=2.
\]
Writing $X=2^{r_a}$, it can be chosen so that
\begin{equation}\label{eq:branchbound}
 R_a<A_a<(X+X^\theta+1)^{1/(a^2-2)}.
\end{equation}
\end{lemma}
\begin{proof}
Put $r_k=C_{k+1}/C_k$. Direct calculation gives
\begin{equation}\label{eq:ratiominus}
 r_k=a+\frac{(a-1)k-1}{a^k-k}>a\ge3.
\end{equation}
Choose $a_0$ so that $2^{r_a}\ge X_0$ for $a\ge a_0$; this is
possible because $r_a\to\infty$. Start with $p_1=2$ and choose
recursively a prime satisfying
\[
 p_k^{r_k}<p_{k+1}\le p_k^{r_k}+(p_k^{r_k})^\theta.
\]
The inputs to the prime interval theorem remain above $X_0$:
$p_k$ strictly increases and $r_k\ge3$, while the initial input is
$2^{r_a}$; more directly, for $k\ge2$, $p_k>2^{r_a}$.
Lemma~\ref{lem:room} gives
\[
 p_{k+1}+1<(p_k+1)^{r_k}.
\]
Consequently the compact intervals
\[
 I_k=[p_k^{1/C_k},(p_k+1)^{1/C_k}]
\]
satisfy $I_{k+1}\subset\operatorname{int}(I_k)$. Their nonempty
nested intersection contains a point $A_a$. For each fixed $k$,
this point belongs to $I_{k+1}$ and therefore to the interior of
$I_k$. It follows that $\floor{A_a^{C_k}}=p_k$ for every $k$.
In particular its first floor is $2$.

At $k=1$, the upper bound on the chosen second prime is
$p_2\le X+X^\theta$. Since $A_a$ lies strictly inside $I_2$,
$A_a<(p_2+1)^{1/(a^2-2)}$, proving the upper bound in
\eqref{eq:branchbound}. The lower bound follows either from strict
containment in $I_1$ or from Proposition~\ref{prop:anchor}.
\end{proof}

\section{The transcendental boundary layer}
\begin{theorem}[Root normalization; \sourceT{57(b)--(c)}]
\label{thm:boundary}
For all sufficiently large integers $a$,
\begin{equation}\label{eq:boundarybound}
 0<\log\frac{\xi_a}{R_a}
 \le\frac{\log(1+2^{-19r_a/40}+2^{-r_a})}{a^2-2},
 \qquad \floor{\xi_a^{a-1}}=2.
\end{equation}
As $a\to\infty$ through integers,
\begin{align}
 \xi_a&=2^{1/(a-1)}
 \left(1+O\left(\frac{2^{-19a/40}}{a^2}\right)\right),
 \label{eq:xiasym}\\
 \Theta_a:=\xi_a^{a-1}
 &=2\left(1+O\left(\frac{2^{-19a/40}}a\right)\right),
 \label{eq:thetaasym}\\
 (a-1)(\xi_a-1)&\longrightarrow\log2.
 \label{eq:limit}
\end{align}
For every integer $a\ge3$, $\Theta_a$ is transcendental and
$\Theta_a>2$.
\end{theorem}
\begin{proof}
By minimality and Lemma~\ref{lem:branch}, $\xi_a\le A_a$ for all
$a\ge a_0$. The branch has $A_a<3^{1/(a-1)}$, while
Proposition~\ref{prop:anchor} gives $\xi_a>2^{1/(a-1)}$.
Thus the minimum also has first floor $2$.

Let $X=2^{r_a}$. Since $X^{1/(a^2-2)}=R_a$,
\eqref{eq:branchbound} implies
\[
 1<\frac{\xi_a}{R_a}
 <(1+X^{-19/40}+X^{-1})^{1/(a^2-2)}.
\]
Taking logarithms proves \eqref{eq:boundarybound} (and a strict upper
bound would also be valid). Because
$r_a=a+1-1/(a-1)$ and $\log(1+u)\le u$ for $u\ge0$,
\[
 \log(\xi_a/R_a)=O(2^{-19a/40}/a^2).
\]
Exponentiating gives \eqref{eq:xiasym}, and multiplying this
logarithmic estimate by $a-1$ gives \eqref{eq:thetaasym}.
Finally,
$2^{1/(a-1)}=1+\log2/(a-1)+O(a^{-2})$, so
\eqref{eq:limit} follows.

Theorem~\ref{thm:family} proves $\xi_a$ transcendental for every
$a\ge3$, not merely for large $a$. A positive integer power of a
transcendental number is transcendental: otherwise it would make the
original number algebraic over $\Qbar$. Thus $\Theta_a$ is
transcendental. Its strict lower bound follows from
Proposition~\ref{prop:anchor}.
\end{proof}

The convergence in \eqref{eq:thetaasym} supplies an explicit family
of transcendental numbers approaching the algebraic value $2$.
It does not suggest that a limit of transcendental numbers must be
transcendental. The threshold $a_0$ comes from the asymptotic prime
interval input and is not numerically specified by this proof.
In particular, the first-floor assertion in
\eqref{eq:boundarybound} has not been proved here for every $a\ge3$.

\section{Exact correction coordinates and finite prefixes}
Define the positive correction coordinate
\begin{equation}\label{eq:tau}
 \tau_a=\log(\xi_a^{a-1}/2).
\end{equation}
Theorem~\ref{thm:boundary} gives
$\tau_a=O(2^{-19a/40}/a)$. At a given exponent $C_k=a^k-k$,
a proposed prime value imposes an exact interval condition.

\begin{proposition}[Prime-cylinder coordinates; \sourceT{57(d)}]
\label{prop:tau}
Let $x=R_a\exp(\tau/(a-1))$, where $\tau\in\R$. For a positive
integer $p$,
\[
 \floor{x^{C_k}}=p
\]
holds if and only if
\begin{equation}\label{eq:tauinterval}
 \frac{a-1}{C_k}\log p-\log2
 \le\tau<
 \frac{a-1}{C_k}\log(p+1)-\log2.
\end{equation}
\end{proposition}
\begin{proof}
The floor equality is equivalent to
$p\le x^{C_k}<p+1$. Substitute the expression for $x$, take
logarithms and multiply by $(a-1)/C_k>0$.
\end{proof}

For a finite prime list $p_1,\ldots,p_m$, simultaneous compatibility
means that the corresponding half-open intervals
\eqref{eq:tauinterval} have a nonempty intersection. There is no need
to rely on rounded root values when certifying their endpoints:
for positive integers $u,v,d,e$,
\[
 u^{1/d}<v^{1/e}\quad\Longleftrightarrow\quad u^e<v^d.
\]
Thus the required comparisons can be reduced to exact integers.
The interval width at exponent $C_k$ equals
\[
 \frac{a-1}{C_k}\log\left(1+\frac1p\right)
 \sim\frac{a-1}{C_kp}.
\]
It becomes small as the output prime and the exponent grow. A tiny
relative correction can therefore change a floor by an integer, and
many more output digits can require many more digits of the constant.

For example, when $a=4$ the first exponents are $3,14,61$.
The anchor $2^{1/3}$ gives
\[
 \floor{(2^{1/3})^{14}}=25=5^2.
\]
Choosing the prime $29$ and putting $x=29^{1/14}$ gives exactly
\[
 \floor{x^3}=2,\qquad\floor{x^{14}}=29,
\]
since $2^{14}<29^3<3^{14}$. This verifies two prime conditions.
It does not identify the global minimum $\xi_4$ or establish that
this particular finite prefix is indefinitely extendible.

\section{What the two-stage construction does and does not give}
A useful unconditional finite-prefix statement follows from the
classical Bertrand theorem. We include its proof to separate it from
the asymptotic use of short prime intervals above.

\begin{proposition}[A two-prime prefix for every $a\ge3$]
\label{prop:two}
Let $X=2^{r_a}$, and let $p$ be the smallest prime strictly greater
than $X$. Then
\[
 y=p^{1/(a^2-2)}
 \quad\text{satisfies}\quad
 \floor{y^{a-1}}=2,\qquad\floor{y^{a^2-2}}=p.
\]
\end{proposition}
\begin{proof}
Let $m=\ceil X$. Bertrand's theorem supplies a prime $q$ with
$m<q<2m$. Hence $X<p\le q<2\ceil X<2(X+1)$.
Since $r_a\ge7/2$,
\[
 (3/2)^{r_a}\ge(3/2)^{7/2}>4
 \quad\Longrightarrow\quad3^{r_a}>4X>2(X+1).
\]
Thus $2^{r_a}<p<3^{r_a}$. Taking $r_a$-th roots gives
$2<y^{a-1}<3$, while the other power equals the prime $p$ itself.
\end{proof}

The proposition begins with a selected and certified prime $p$.
It records how that prime is encoded in a two-stage constant.
It does not produce a prime without finding one, does not give the
exact global minimum, and does not make the uncorrected value
$\floor{2^{r_a}}$ prime. For sufficiently large $a$, the short-interval
argument can instead choose the smallest prime above $X$ inside the
Baker--Harman--Pintz interval and extend that prefix through nested
cylinders. The numerical threshold remains unspecified.

\section{Scope and source correspondence}
Theorems~\ref{thm:residues} and \ref{thm:composite} retain the
source content T47 and T49, with the coefficient extension from T56.
Theorem~\ref{thm:family} combines T48 and T56(a), so the
$b=1$ special case is not counted as an additional independent
method. Theorem~\ref{thm:braids} and
Propositions~\ref{prop:scaling}--\ref{prop:sync} contain the remaining
structural portions of T56. Sections~9--13 establish T57 with its
large-$a$ quantifiers explicit. The separate source item T53,
concerning $\xi(2^k)$, is not part of this manuscript.

The general congruence lemma applies already at $a=2$, but the
transcendence application in Theorem~\ref{thm:family} is stated only
for $a\ge3$. The limit of consecutive ratios must exceed $40/19$
to invoke the external criterion used here. An endpoint argument for
$a=2$ would require additional input.

The set
\[
 \{x>1:\floor{x^{3^k-k}}\text{ and }\floor{x^{3^k+k}}
 \text{ are prime for every }k\ge1\}
\]
is not proved nonempty here. The two separate nonempty sets need not
intersect. Interleaving their exponent sequences gives adjacent ratios
approaching $1$ within each pair, so the Type-B growth assumptions
cannot simply be inherited. Likewise, no monotonicity of $\xi(C)$
under pointwise comparison of different exponent sequences has been
proved.

The logical contribution to Problem~1.7 is an explicit family of
intrinsically non-reversible recurrences with universal congruence
visitation, combined with an established composite-progression input.
The transcendence contribution is verification of a separate,
established criterion for new explicit exponent families and their
stated braid variants. Neither contribution settles the transcendence
of the classical Mills constant, and neither is presented as a proof
of historical firstness. All historical priority assertions would
require an independent literature audit beyond the proof checks
recorded here.

\begin{thebibliography}{9}
\bibitem{SaitoIntervals}
K.~Saito, \emph{Intervals without primes near an iterated linear
recurrence sequence}, arXiv:2504.14968, version~1, 21 April 2025,
Problem~1.7 and the introductory statement (D1).
\href{https://arxiv.org/abs/2504.14968}{arXiv:2504.14968}.
Published in \emph{International Journal of Number Theory} (2026),
\href{https://doi.org/10.1142/S1793042126500685}
{doi:10.1142/S1793042126500685}.
\bibitem{SaitoVariants}
K.~Saito, \emph{Transcendency of variants of Mills' constant},
arXiv:2508.16068, version~3, 7 December 2025, Theorems~1.6 and~2.3.
\href{https://arxiv.org/abs/2508.16068v3}{arXiv:2508.16068v3}.
\bibitem{Dubickas}
A.~Dubickas, \emph{Intervals without primes near elements of linear
recurrence sequences}, International Journal of Number Theory
\textbf{14} (2018), no.~2, 567--579.
\href{https://doi.org/10.1142/S1793042118500355}
{doi:10.1142/S1793042118500355}.
\bibitem{BHP}
R.~C.~Baker, G.~Harman and J.~Pintz,
\emph{The difference between consecutive primes, II},
Proceedings of the London Mathematical Society (3)
\textbf{83} (2001), no.~3, 532--562.
\href{https://doi.org/10.1112/plms/83.3.532}
{doi:10.1112/plms/83.3.532}.
\end{thebibliography}
\end{document}
